\chapter{When a Collective Can Act as an Agent}
\label{comp:inheritance}

A collective can have a simple description without being able to evolve from
that description alone. Knowing the total energy of two bodies, for example,
does not usually determine their next interaction. A parent of law-valued
agents faces the same question: does its retained state determine its future
responses?

There are two operations to distinguish. A channel $C$ compresses represented
contents and gives a parent law $Q_xC$. A map $R$ retains a complete dynamical
state $R(x)$. Besides belief and model laws, that state can include the model
evaluator, source information, boundary channels, and memory. A distribution
over possible contents does not automatically contain these additional data.

\section{Two routes to the same parent}

Suppose the fine configuration advances by $F$. Exact inherited evolution
requires a parent update $F^c$ such that
\begin{equation}
 R(F(x))=F^c(R(x)).
\end{equation}
We can first evolve the children and then form their parent, or first form
the parent and evolve it. Both routes must give the same retained state.
Repeating the identity gives agreement at every later step. Boundary response
mechanisms must likewise depend only on the retained state.

\begin{figure}[t]
\centering
\begin{tikzpicture}[
  >={Stealth[length=2.2mm]},
  box/.style={draw=primaryblue,rounded corners=2pt,fill=primaryblue!5,minimum width=2.55cm,minimum height=0.95cm,align=center,font=\small},
  every node/.style={font=\small}]
\node[box] (fine) {Fine state\\$x$};
\node[box,right=4.3cm of fine] (finenext) {Fine successor\\$F(x)$};
\node[box,below=1.8cm of fine,fill=secondaryblue!10] (parent) {Parent state\\$R(x)$};
\node[box,below=1.8cm of finenext,fill=secondaryblue!10] (parentnext) {Parent successor\\$R(F(x))=F^c(R(x))$};
\draw[->,thick,primaryblue] (fine) -- node[above=4pt] {fine update $F$} (finenext);
\draw[->,thick,secondaryblue] (parent) -- node[below=4pt] {inherited update $F^c$} (parentnext);
\draw[->,thick,primaryblue] (fine) -- node[left=5pt] {$R$} (parent);
\draw[->,thick,primaryblue] (finenext) -- node[right=5pt] {$R$} (parentnext);
\node at ($(fine)!0.5!(parentnext)$) {same retained result};
\end{tikzpicture}

\caption{The two routes to the next parent configuration must agree.
Vertical maps retain information; horizontal maps advance the dynamics.
A compressed law alone need not make this square commute.}
\label{fig:comp-inheritance}
\end{figure}

The immediate test is whether configurations with the same parent have the
same next parent. If $R(x)=R(x')$ but $R(F(x))\ne R(F(x'))$, no single
$F^c$ can satisfy the equation. Conversely, in a finite state space, this
test is sufficient: evolve any representative of each parent state.
Agreement within each fiber of $R$ makes the choice irrelevant. On general
measurable spaces the resulting map must also be measurable.

For stochastic dynamics, replace the next state by its distribution. If $T$
is the fine transition matrix, the probability
\begin{equation}
 \sum_{y:R(y)=z'}T(x,y)
\end{equation}
must depend only on $R(x)$. These sums then define normalized parent rows.
Conditioning on any parent history merely averages identical rows, so the
parent reproduces the whole retained history law. This is strong lumpability
\citep{KemenySnell1976}. Agreement for one equilibrium preparation is a
weaker statement.

\section{When a gradient flow descends}

Suppose the fine configuration lies on a Riemannian manifold with metric $g$ and evolves by $-\operatorname{grad}_g E$. Let a smooth surjective submersion $R$ retain the parent, whose metric is $g_c$. One sufficient condition is that the energy depend only on the parent,
\begin{equation}
E=E_c\circ R,
\end{equation}
and that $R$ preserve lengths of horizontal vectors, meaning vectors orthogonal to $\ker DR$. Then the gradient is horizontal because the energy has no derivative along a fiber of $R$. For a horizontal $u$,
\begin{equation}
g_c(DR\operatorname{grad}_gE,DR u)
=g(\operatorname{grad}_gE,u)
=dE_c(DR u).
\end{equation}
Every parent tangent has a horizontal lift, so the pushed gradient is exactly $\operatorname{grad}_{g_c}E_c$. The dynamics descends in the same time.

If instead $E=\chi(E_c\circ R)$ and the horizontal metric is scaled by $d(x)^2$, the pushed vector field acquires the factor $\chi'(E_c(Rx))d(x)^2$. A positive factor preserves oriented trajectories after a time change. A factor that depends on discarded state does not specify an autonomous parent clock. Taking the minimum of $E$ over each fiber does not supply the first hypothesis: it gives the value at the minimizing states, not the fine energy everywhere on that fiber.

The requirement $E=E_c\circ R$ is stronger than descent needs, and it has a physical cost: it makes the energy blind to every motion along a fiber, so the discarded coordinates cannot move relative to the retained ones. Only the horizontal derivatives matter. The computation above uses $dE(u)$ for horizontal $u$ alone, so descent holds exactly when
\begin{equation}
dE(u)=dE_c(DR(u))\qquad\text{for every horizontal }u.
\end{equation}
Two particles on a line with unit masses, harmonic self-energies $\tfrac12x_1^2+\tfrac12x_2^2$, and any pair interaction $W(x_1-x_2)$ show the difference. The center of mass $(x_1+x_2)/\sqrt2$ relaxes as if the interaction were absent, because a force that depends only on the separation pushes the two particles in opposite directions and cancels in their sum. Meanwhile the separation evolves under $W$. The parent is exact, the discarded coordinate has its own dynamics, and the energy is not a function of the center of mass, except for the single interaction $W(s)=-s^2/4$ (up to a constant) that exactly cancels the separation part of the harmonic self-energies and freezes the separation.

The opposite case is a retention that simply drops some agents. If the retained and discarded agents exchange any peer cost, the retained velocity depends on the discarded state, because each agent pulls on every agent it attends to and is pulled back. Dropping whole agents is exact only for blocks that do not interact at all. Exact compression of interacting agents can come from symmetry, as when every member of a block can be rotated together without changing any cost, or from special energy structure such as the center-of-mass cancellation; these are two sufficient routes, and the general test is whether the retained velocity is constant along each fiber. A related consequence is that exact inheritance creates no oscillation at the image of an equilibrium of the fine flow: the linearization of a gradient flow at a critical point has real eigenvalues, and an exact parent inherits a subset of them. A parent equilibrium need not be such an image. If one agent drifts steadily while never reaching a critical point, and the parent reads a second, relaxing agent in a frame turned by the first, the parent spirals into its equilibrium although every fine motion is a gradient descent.

\section{A missing role can prevent closure}

Let two independent represented bits have laws
$\operatorname{Ber}(p_1)$ and $\operatorname{Ber}(p_2)$. Their probabilities
evolve according to
\begin{equation}
 \begin{aligned}
 \dot p_1&=-p_1+\lambda(p_2-p_1),\\
 \dot p_2&=-2p_2+\lambda(p_1-p_2),\qquad \lambda\geq0.
 \end{aligned}
\end{equation}
The square $[0,1]^2$ is preserved: the vector field never points outward
on its boundary. It is tangent to a face when $\lambda=0$ and vanishes at the
corner $(0,0)$. Retain only the law of the sum of the two bits. The configurations
$(0.8,0.2)$ and $(0.2,0.8)$ both give probabilities
$(0.16,0.68,0.16)$ for sums $0,1,2$.

Their subsequent changes differ. The coupling terms cancel when we add the
two equations, giving
\begin{equation}
 \frac{d}{dt}\E(K_1+K_2)=-(p_1+2p_2).
\end{equation}
The derivatives are $-1.2$ and $-1.8$. Thus the same retained law has two
different possible derivatives. Compression erased which probability belonged
to the faster-decaying role. The fine recognition law remains a product
throughout. What must be retained is the information this evolution uses;
cross-agent covariance is not automatically required.

\section{An exact supplied hierarchy}

A finite construction shows that two successive reductions can work. Take
$x=(a,b,c,r)\in\{0,1\}^4$, with internal register $r$ and readout
$\rho(x)=r$. Let the agent laws be
\begin{equation}
 q_a=\operatorname{Ber}\left(\frac{1+2a}{4}\right),\qquad
 s_a=\operatorname{Ber}\left(\frac{2+a}{4}\right).
\end{equation}
The model evaluator assigns probability $1/4+m/2$ to record value one
when the model presentation is $m\in\{0,1\}$. Averaging over $s_a$
gives the predictive law $\nu_a=\operatorname{Ber}(1/2+a/8)$.

At each step draw independent bits $\xi,\eta,\upsilon$ with probabilities
$1/4,1/3,1/5$ of being one, independently of previous draws, and set
\begin{equation}
 \begin{aligned}
 a'&=a\mathbin\oplus\xi,\\
 b'&=b\mathbin\oplus a\mathbin\oplus\eta,\\
 c'&=c\mathbin\oplus b\mathbin\oplus\upsilon.
 \end{aligned}
\end{equation}
Here $\oplus$ means addition modulo two. Independently of these draws,
conditional on the current state, give $r'$ law $\nu_a$. The record is
therefore an internal state change. Its generation rule is supplied, rather
than deduced from the existence of an evaluator.

First retain $(a,b,r)$, then retain $(a,r)$. Summing the fine transition
over $c'$ removes a normalized factor and leaves a row independent of $c$.
Summing that row over $b'$ removes another normalized factor and leaves a
row independent of $b$. The configuration spaces shrink from sixteen to
eight to four states. Both parents retain $q_a,s_a$, their evaluator, and
the register. They reproduce the complete retained state and record
histories for every initial law. The coupling is real, since $a$ affects
$b'$ and $b$ affects $c'$; its triangular form permits these reductions.
This is a drive and its response: $a$ drives $b$, which drives $c$, and
nothing acts back. Each retained block is autonomous because the variables
it keeps never read the variables it discards. The configuration describes
one agent with internal marks, so neither level groups two agents, and the
hierarchy itself was prescribed.

\section{What a parent must remember}

When closure fails, a parent can retain more state or, under a declared
Markov process, the conditional law of the current hidden state given its
accessible history. The Markov property matters: if the next state can depend
on older states directly, two histories with the same current conditional law
can predict different futures. In a finite model this is a probability
vector: propagate it with
the fine transition, restrict to the new retained value, and normalize.
That vector determines every future response even when the current summary
does not. Keeping only an equilibrium hidden distribution generally changes
transient predictions. The vector lives on the simplex of all fine states,
so it saves nothing unless noise or poor observability keeps the history from
pinning down the fine state.

There is a middle case in which the parent needs no memory although the
current summary is not autonomous for every preparation. Suppose that for
each retained value $z$ there is a fixed way $\Lambda(z,\cdot)$ of spreading
probability over the fine states with that value, and that spreading and then
evolving gives the same result as evolving the parent and then spreading,
$\Lambda T=T^c\Lambda$. A child population prepared by spreading a parent law
in this way stays spread in this way forever: the hidden distribution is
always $\Lambda(Z_n,\cdot)$, and the parent is an ordinary Markov chain with
kernel $T^c$ \citep{RogersPitman1981}. A handoff that prepares children
differently can break the parent, even though the kernels are unchanged.

Given the fine update and a list of what the parent must reproduce (its
laws, the responses to a declared set of probes, and the channels it reads),
there is a smallest exact parent. For deterministic dynamics two
configurations are lumped together exactly when the listed quantities agree
for them now and at every later time. Nothing coarser can be exact, and
anything exact remembers at least this much. For random dynamics the
smallest exact parent is finer: two states are lumped when they send equal
probability, jointly with each listed response, into every lumped class,
which can separate states whose future listed quantities have the same
laws. A longer list of probes gives a finer parent; if outside agents
may probe the children without restriction, they can tell every exposed child
law apart and the parent cannot forget any of them. How much a parent
forgets is therefore decided by the probes declared before the parent is
chosen. ``Before'' has a definite meaning here: the probes are given by a
rule that assigns a class to every possible block from the constitutive data
alone. A list written down after the retention, such as ``every probe whose
answer the parent can reproduce'', would make any retention exact.

There is an exact way to separate insufficient memory from inaccurate prediction. Let $P$ be a declared law of a finite record sequence $Z_0,\ldots,Z_{T-1}$ on standard Borel spaces. A retained statistic is $M_t=f_t(Z_{<t})$. Let $S$ use normalized prediction rows $s_t(\cdot\mid M_t)$. The path-law chain rule and conditional disintegration give
\begin{equation}
\begin{split}
\KL(P\|S)
={}&\sum_{t=0}^{T-1}I_P(Z_t;Z_{<t}\mid M_t)\\
&+\sum_{t=0}^{T-1}\E_P
\KL\bigl(P(Z_t\in\cdot\mid M_t)\|s_t(\cdot\mid M_t)\bigr).
\end{split}
\end{equation}
The first sum is information about the next record that was discarded with the history. The second is error in the chosen prediction given the retained information. Both are nonnegative, including when the total is infinite. This extends the finite-window Markov decomposition \citep[Thm.~7.3.1]{Gray1990}. Zero first sum gives predictive sufficiency for the declared process; a compact recursively updated memory still needs a realization $M_{t+1}=U(M_t,Z_t)$. Changing the allowed controls changes which histories must be predicted.

Statistical dependence alone also does not select an agent hierarchy. A declared finite cross-agent joint law has a finest partition into independent blocks, but arbitrarily weak couplings can merge those blocks. A deterministic selector reading only a fully permutation-symmetric law cannot choose a nontrivial intermediate hard partition while respecting relabeling: swapping one member with an outside agent would create a different block that overlaps the original. More generally, such a selector can only output a partition that every symmetry of the law preserves; on a ring of four agents with the symmetries of a square, it can pair opposite agents but never neighbors. Realized histories, additional geometry, or a distribution over partitions can fall outside that restriction. The result identifies what a selector must read; it does not identify independence with agency.

\section{The formation standard}

The theory aims to show that interacting agents can produce further agents.
A construction that groups agents and hands their laws to parents does not
yet show this, because the grouping and the handoff can be arranged without
the interaction doing anything. The formation standard of the main text
lists six clauses, each closing one such loophole, and it asks that every
choice that could otherwise be made after seeing the dynamics be fixed
first: the controls, the sense of ``open'', the probes, the selector and its
constants, the time window, and the counterfactual handoffs.

The first clause asks for a set of starting configurations that is open in
the coordinates in which the dynamics is actually posed, or open inside a
family preserved by the dynamics and defined by stated constraints, with
that weaker grade reported. A single trajectory is not enough. Along one
path, almost any summary of the state traces out a curve, and every curve is
the solution of some differential equation, so ``the parent evolves
autonomously'' would hold for nearly any choice of parent. Exactness means
something only when it is tested on configurations that the parent cannot
tell apart. Declaring an arbitrary topology would not help, since a single
point is open in the discrete topology. The second clause asks the dynamics
to pick nontrivial groups that persist by a stated criterion, with a
selector drawn from a family whose constants are fixed in advance and range
over an interval rather than a single tuned value. The third asks that each
group carry inherited laws, evaluators, responses, and dynamics for the
probes that outsiders may use, given by the rule described in the previous
section.

The fourth clause is compression. A parent that simply keeps every child's
state is always exact and says nothing. Four agents with Gaussian laws pulled
together by symmetric peer costs make the point: a threshold rule first pairs
the two closest agents, then the other two, then the two pairs, and keeping
the full tuple of child states at each stage reproduces the dynamics exactly.
Nothing has been forgotten, so nothing has formed. The clause therefore asks
each parent to lump together configurations that a child-level probe could
still distinguish. These probes are what the agents, their channels and
records, and any declared external querier can do to individual children
before the handoff. A distinction that none of them can resolve is gauge in
the sense of the first chapter, and erasing it is not compression.

The fifth clause is the control experiment. The main text declares two
nulls. In the label-preserving null every peer cost is deleted while the
sources keep their labels and prior weights; a receiver that also compares
itself with its own model then gives that self comparison a weight below
one, because the deleted peers still hold part of its row at zero cost. In
the renormalized null the deleted peers are removed from the row, so the self
comparison keeps weight one. When an agent's self cost is nonzero the two
nulls change the self force differently, and a grouping difference that
appears against only one of them may come from that reweighting rather than
from the interaction. The standard therefore asks for a difference against
both. Only costs that actually read another agent's law are deleted. A cost
whose comparison law ignores every other agent, such as one through a
channel with a fixed output, stays in place: deleting it would move each
agent's own target, and agents that never interact could then pass the
control. A third control, the frozen environment, deletes nothing but holds
every other agent's law at its starting value in each agent's force. It asks
a different question, whether each agent's response to where the others
started already produces the groups, without any feedback.

The standard asks that the sequence of selected groups differ from the null
one at every point of an open set of starting configurations. Only the
groups count. If the interaction merely makes the same joins happen sooner
or later, or reverses the order of two joins of disjoint groups, it has
changed a clock, not the outcome. On a finite time window the comparison has
to be made with care: if the interaction doubles the speed of a common
clock, the full flow can complete within the window joins that the null
completes only later. The standard therefore compares the full sequence on
its window with the beginning of the complete null sequence. Agents that do
not interact can be grouped by any rule that reads their initial data and a
common relaxation clock, much as runners in a race can be grouped by
finishing order even though they never touch, and the standard excludes
exactly that.

The verdict of the fifth clause is reported with the selector constants for
which it holds and with a quantity that does not depend on the selector at
all. Without that discipline a threshold placed after the trajectories are
known can decide the verdict. A stronger grade asks for a difference at every
starting point and for every constant, and an amplifying grade asks that
groups be tighter relative to their separations than the same agents, on
the same partition, at the same times, without interaction.

A whole class of architectures fails this control. If every agent has the
same model-independent evaluator, no agent attends to itself, and every
channel delivers one common law that depends on the population only through
an additive moment, then each pair of agents relaxes toward each other at a
common rate: the logarithm of the ratio of their laws shrinks by the factor
$e^{-2t}$, up to a constant. Without the peer costs the same quantity shrinks
by $e^{-t}$. More precisely, the full flow is the null flow run at doubled
speed and then tilted by one common factor. Any selector that is blind to
such a common tilt, for example one that compares pairs through their law
ratios, produces the same sequence of groups with or without the interaction,
at different times. A selector that is not blind to it, for example a
threshold on the divergence between two laws, can give different groups,
which is why the selector is fixed in advance and a selector-free quantity
is reported. Once agents also compare themselves with their own models
through the common channel, the two nulls no longer share a clock, and the
argument does not apply. A related statement holds for scalar parameters
moved by a common state: if summing two agents' parameters is an exact parent
for every pair, the parameter equation is affine, all differences shrink by
a common factor, and the order of joins along the full flow is fixed by the
initial distances. That alone does not decide the control, which also
depends on how the null flow moves.

The sixth clause is recursion. The parents must themselves be grouped, by
the inherited interactions and the same selection rule, without resetting
their data or adding a new attraction tuned for the purpose, and this second
grouping must depend on what the first stage produced. The test uses a set
of counterfactual handoffs fixed in advance: the actual handoffs, together
with handoffs made at the same times from runs in which the first groups had
not yet formed. If the second grouping happens from all of them except a set too thin to contain an open piece of the class, or from every one of
them, it is a sequence of scales supplied in the initial data, not a parent
built from parents. Counterfactual parents that already sit inside the
joining threshold are left out of this set, since a rule that joins groups
when their distance falls through a threshold never joins them, and
including them would let every such construction pass.

A construction that meets the selection, inheritance, and first recursion
requirement but is shown to fail compression, the control, or the
counterfactual test is called a two-stage selection with exact inheritance.
When a clause is neither established nor shown to fail, its status is left
open and named. The standard does not ask that a parent be an agent of the
same type as its children or that it act back on them, so two stages here do
not yet mean a hierarchy that repeats itself.

The main text calibrates the clauses on four Gaussian agents pulled toward
their own fixed anchors and toward each other. The coupling shrinks every
difference to a third of the corresponding anchor difference, so two pairs
form and then the pair of pairs, while without the coupling nothing forms at
all; the control is passed, even in its strong grade. The frozen
environment already reproduces the groups. And the union of the two parents
forms from every counterfactual handoff, because the gap between two parents
always relaxes to the same value set by the anchors. The example is
therefore a two-stage selection that fails the counterfactual test, which is
what it should be: its stages are supplied by the anchors.

The constructions in the later chapters are measured against this standard,
each with a report of the same form. The nonsmooth fusion example produces
identical fused children, so it compresses nothing. The common-precision
matrix-law construction of \cref{cf:formation} selects two parents and then
their union with a schedule that does not depend on the coupling; there the
two nulls differ, because each agent compares itself with its own model, and
the control passes against one and fails against the other, so it fails. The
six-agent construction of \cref{cf:open-formation} works on a full open
neighborhood of laws with selective channels, but its theorem does not
settle the control or the counterfactual test, and integrations of a limit
field indicate that the null flow forms the same groups. Both smooth
constructions supply their compression architectures in advance. No
construction in this manuscript is proved to meet every clause.
